\section{Related Work}
\label{sec:related}

Four lines of work bear on our questions. \Cref{tab:related} places them side by side.

\paragraph{Cocoons measured over categories.} The idea that personalised media can enclose people goes back to Sunstein's information cocoons~\citep{sunstein2006infotopia,sunstein2017republic} and Pariser's filter bubble~\citep{pariser2011filter}. \citet{michiels2022filterbubbles} reviewed the literature, found the definitions in use mutually inconsistent, and proposed a testable one: a decrease over time in the diversity of what a user is recommended. Empirical studies have followed this reading and measured diversity over content categories. On MovieLens, consumption narrowed slightly over time, and less for users who followed recommendations~\citep{nguyen2014exploring}; on Spotify, algorithmic listening went with lower diversity~\citep{anderson2020spotify}. On Kuaishou, \citet{li2022cocoon} compared the category diversity of what new users were shown and what they watched, and \citet{sukiennik2024deep} showed narrowing inside a broad interest, for instance from sport to one club. \citet{piao2023cocoons} explained such narrowing by the interplay of similarity-based matching and feedback. All of these treat a cocoon as concentration over topics. A user who watches many marriage videos with a balanced range of views scores as more ``cocooned'' than one who watches few and sees mostly a single hostile narrative.

\paragraph{Audits at the level of viewpoints.} A second line measures stance or ideology, almost always for politics or misinformation. With about 100{,}000 automated accounts, \citet{haroon2023youtube} found YouTube's recommendations mostly congenial to an account's ideology; \citet{hosseinmardi2024bots} found that the sidebar forgets a partisan history within about 30 videos once the account switches to moderate content. Watching debunking videos can burst misinformation bubbles~\citep{tomlein2021audit,srba2023auditing}, and on TikTok the feed reinforces an account's stance on polarising topics~\citep{pecher2026tiktok}. These audits see the platform from outside and so cannot say whether a skew comes from what is available, from ranking that affects everyone, or from personalisation. Studies with platform data can: \citet{bakshy2015exposure} separated the roles of the friend network, ranking and choice on Facebook, and \citet{huszar2022amplification} measured how much Twitter's ranking amplifies political content relative to a chronological feed. Large experiments have since found that changing the feed alters exposure far more than attitudes~\citep{guess2023feed,nyhan2023likeminded}, which is why we claim results about exposure only. None of this work covers the everyday social issues of short-video feeds, or a Chinese platform.

\paragraph{What user controls achieve.} In donated data from 22{,}722 YouTube users, ``not interested'' prevented about 11\% of unwanted recommendations and ``don't recommend channel'' about 43\%~\citep{ricks2022button}; with a different design, \citet{liu2024train} found that ``not interested'' removed most recommendations from an unwanted topic. On TikTok the same button was the most effective explicit signal, yet many feeds reverted once accounts stopped pressing it~\citep{kaplan2026foryou}. \citet{vombatkere2024tiktok} measured how a feed balances exploring and exploiting interests, and recommender designs have been proposed that let users ask for less of a category~\citep{wang2022controllable}, counteract over-exposure~\citep{gao2024cirs}, or learn from negative feedback~\citep{wang2023negative}. These studies work at the level of topics, channels or similar videos, and report what is removed, not what is lost with it. They do not ask whether a control changes the mix of views within a topic the user wants to keep.

\paragraph{Feedback loops in models and simulation.} Theory shows that recommendation feedback can homogenise behaviour~\citep{chaney2018confounding}, drive interests to degenerate states~\citep{jiang2019degenerate} and amplify preferences~\citep{kalimeris2021amplification}; \citet{dean2020reachability} ask which items a user can make reachable by changing their ratings. Recent simulations replace rule-based users with language-model agents~\citep{zhang2024agent4rec,sukiennik2025simtok,yang2026recloop}, and \citet{feng2026bep} quantify the potential to escape a bubble by contrasting simulated behaviours. In these models the effect of each action is set by hand or produced by an agent that is hard to validate, the recommender is rarely one trained on a platform's own logs, and the outcome is again category diversity.

\paragraph{The gap.} Prior work measures cocoons over categories, audits viewpoints for political content from outside the platform, evaluates controls by what they remove, and simulates feedback with assumed responses. We measure skew \emph{within} an everyday issue separately from interest in it, attribute it to supply, platform and personalisation, read the recommender's response off a model trained on the platform's logs, and price each escape route in lost topic, effort and persistence. The analysis of engagement connects to the finding that engagement-based ranking amplifies divisive content~\citep{milli2025engagement}, and the design in \cref{sec:design} to earlier tools that show users the balance of their reading~\citep{munson2013balance,nagulendra2014visualization}.

\begin{table}[t]
\caption{Prior work and this study. ``Separates layers'' asks whether a skew is attributed to supply, platform-wide ranking and personalisation; ``topic cost'' whether the evaluation of a user action reports how much of the wanted topic is lost.}
\label{tab:related}
\centering\small
\setlength{\tabcolsep}{4.5pt}
\begin{tabularx}{\textwidth}{@{}L{3.5cm} Y C{1.2cm} C{1.4cm} C{1.1cm} C{1.3cm}@{}}
\toprule
Line of work & What is measured & Within-topic view & Separates layers & Topic cost & Detect\-ability\\
\midrule
Category-level cocoons \citep{nguyen2014exploring,li2022cocoon,sukiennik2024deep,anderson2020spotify} & diversity of consumption or exposure over categories & \xmark & \xmark & \xmark & \xmark\\\addlinespace[2pt]
Outside audits \citep{haroon2023youtube,hosseinmardi2024bots,srba2023auditing,pecher2026tiktok} & stance of recommendations to scripted accounts (politics, misinformation) & \cmark & \xmark & \xmark & \xmark\\\addlinespace[2pt]
Platform-data studies \citep{bakshy2015exposure,huszar2022amplification,nyhan2023likeminded} & ideological exposure and amplification with internal data & \cmark & \pmark & \xmark & \xmark\\\addlinespace[2pt]
User-control studies \citep{ricks2022button,liu2024train,kaplan2026foryou} & share of unwanted topic or channel removed by a control & \xmark & \xmark & \xmark & \xmark\\\addlinespace[2pt]
Models and simulations \citep{chaney2018confounding,jiang2019degenerate,dean2020reachability,feng2026bep} & homogeneity, degeneracy, reachability, category diversity & \xmark & \xmark & \xmark & \xmark\\\addlinespace[2pt]
\rowcolor{rowhl}\textbf{This study} & skew of a narrative within an everyday issue & \cmark & \cmark & \cmark & \cmark\\
\bottomrule
\end{tabularx}
\end{table}
